\begin{center}
    \Large
    \textbf{Embracing Respiration-Induced Motion in Cardiovascular MRI}

    \vspace{0.4cm}
    \normalsize
    Tarun Naren
    
    \vspace{0.4cm}
    \normalsize
    \textit{Under the supervision of Dr.\ Oliver Wieben \\
    at the University of Wisconsin-Madison}
    
    \vspace{0.9cm}
    \Large
    \textbf{Abstract}
\end{center}

\doublespace{}

Respiratory motion is a well-known complicating factor in magnetic resonance imaging (MRI), especially when imaging elements of the cardiovascular system. Periodic motion of the lungs during the respiratory cycle affects the torso and abdomen, which impacts imaging of the heart, great vessels, and the liver. The resulting motion artifacts, which appear in traditional Cartesian imaging as distinct replicas in the phase-encoding direction, negatively impact the quality of diagnostic information that can be gleaned from MR images. One common way to mitigate these effects is to instruct subjects to hold their breath for the duration of the acquisition. However, breath holds can be infeasible for long acquisitions or for subjects that have difficulty holding their breaths. Additionally, respiration has important physiological consequences and there is merit to exploring the impacts of respiration and respiration-induced motion on the cardiovascular system.

In this thesis, I seek to implement and optimize methods for not only mitigating the effects of respiration in cardiovascular MRI which would enable otherwise impossible imaging, but to also utilize information from the respiratory cycle to extract valuable diagnostic information. This will be performed in three different scientific contexts, each focusing on a different anatomical region: 1.\ improving radial, free-breathing 2D Phase Contrast (2D PC) MRI for measuring pulse wave velocity in the aorta, 2.\ accelerating the acquisition of radial, free-breathing 4D flow MRI in the portal venous system, and 3.\ implementing a 5D MRI approach for resolving cardiac and respiratory motion of the heart in radiation therapy.